\section{Validación experimental}
\label{sec:validacion}

La validación se organiza en los tres niveles que corresponden a un trabajo de instrumentación: la caracterización electrónica del instrumento, una prueba controlada de adquisición y una campaña de campo representativa. El caso de campo se presenta como demostración de la utilidad del instrumento y no como un estudio geofísico autónomo.

\subsection{Caracterización electrónica del frente analógico}
\label{sec:ensayo-electronico}

El ensayo buscó evaluar si la respuesta del compensador, combinada con la del geófono, permite extender el ancho de banda del sensor hacia frecuencias inferiores a su resonancia. Se inyectaron barridos senoidales de tensión y se registraron simultáneamente, con cuatro canales de osciloscopio, las salidas de la ganancia programable, la rama pasabanda, el compensador y el filtro pasa-bajos. A partir de las relaciones entre canales se estimaron la magnitud, la fase y la coherencia de cada transferencia. Los barridos cubrieron bandas superpuestas desde 10 mHz hasta 200 kHz.

El estimador empleado opera con las señales analíticas de entrada $u_a[n]$ y salida $y_a[n]$, obtenidas mediante la transformada de Hilbert. Para las muestras $\mathcal B_k$ del intervalo de frecuencia $k$:
\begin{equation}
\begin{aligned}
C_k&=\sum_{n\in\mathcal B_k}y_a[n]u_a^*[n],\\
E_{u,k}&=\sum_{n\in\mathcal B_k}|u_a[n]|^2,\qquad
E_{y,k}=\sum_{n\in\mathcal B_k}|y_a[n]|^2,\\
\widehat H_k&=\frac{C_k}{E_{u,k}},\qquad
\gamma_k^2=\frac{|C_k|^2}{E_{u,k}E_{y,k}}.
\end{aligned}
\label{eq:estimacion-transferencia}
\end{equation}
El asterisco indica conjugación compleja, $\widehat H_k$ es la transferencia estimada y $\gamma_k^2$ la coherencia. Se excluyen intervalos saturados o con excitación insuficiente.

La respuesta conjunta sensor--acondicionamiento se calculó combinando el modelo nominal del geófono con la transferencia electrónica medida. Un ensayo con el sensor sometido a una excitación mecánica conocida permitiría verificar directamente esa respuesta, pero no se dispuso del equipo necesario. La inyección de tensión permitió caracterizar por separado el circuito y evaluar su contribución a la extensión de banda.

Para comparar las mediciones con el diseño se consideraron dos intervalos: el generado por las tolerancias de los componentes, manteniendo el potenciómetro en el ajuste nominal, y el alcanzable al recorrer el potenciómetro. La tabla~\ref{tab:envolventes} presenta la proporción de puntos comprendidos en cada intervalo.

\begin{table*}[!ht]
\centering
\footnotesize
\caption{Identificación de las etapas y ubicación de la medición respecto de
dos envolventes: la de tolerancias de componente con el potenciómetro en su valor de
diseño, y la que barre todo el recorrido del potenciómetro. \revsolution{Fuente: elaboración propia.}}
\label{tab:envolventes}
\begin{tabularx}{0.92\textwidth}{YYYYYY}
\toprule
Ruta medida & Puntos & Banda identificada & Coherencia & Dentro de tolerancias (mag./fase) & Dentro del recorrido del pot.\ (mag./fase) \\
\midrule
BP & 85 & 0,21 Hz – 47,5 kHz & 0,9684 & 32 \% / 6 \% & 32 \% / 6 \% \\
Compensador & 71 & 0,21 Hz – 4,74 kHz & 0,9643 & 37 \% / 11 \% & \textbf{92 \% / 94 \%} \\
LP & 33 & 11,75 Hz – 1,11 kHz & 0,9454 & 82 \% / 94 \% & 82 \% / 94 \% \\
\textbf{PGA → ADC (ruta completa)} & \textbf{61} & \textbf{0,21 Hz – 1,11 kHz} & \textbf{0,9946} & 41 \% / 15 \% & \textbf{100 \% / 100 \%} \\
\bottomrule
\end{tabularx}
\end{table*}

En la ruta completa PGA--ADC, el 100 \% de los puntos queda dentro del intervalo alcanzable con el potenciómetro, tanto en magnitud como en fase; para el compensador, las proporciones son del 92 \% y del 94 \%. La respuesta de la cadena es, por tanto, compatible con el rango de ajuste previsto. Las tolerancias por sí solas, con el potenciómetro fijo en su valor nominal, no explican toda la diferencia: la posición del ajuste es determinante.

\begin{figure*}[!tb]
\centering
\includegraphics[width=\textwidth,keepaspectratio]{f_ruta_pga_adc.png}
\caption{Identificación de la cadena completa PGA--ADC. De arriba abajo: magnitud, fase y coherencia local. Se comparan la medición cruda, la procesada, el modelo identificado y el modelo nominal del circuito con los componentes actuales. \revsolution{Fuente: elaboración propia a partir de los registros del proyecto.}}
\label{fig:identificacion}
\end{figure*}

\begin{figure*}[!tb]
\centering
\includegraphics[width=\textwidth,keepaspectratio]{f_identificacion_comp.png}
\caption{Identificación de la rama de compensación respecto de la etapa de ganancia. El apartamiento respecto del modelo nominal se explica por la posición del potenciómetro y no por un componente fuera de especificación. \revsolution{Fuente: elaboración propia a partir de los registros del proyecto.}}
\label{fig:identificacioncomp}
\end{figure*}

El aprendizaje principal fue la sensibilidad de la cancelación al valor del potenciómetro. Con los componentes nominales, la cancelación exacta y el ajuste objetivo están separados por apenas 1,98 $\Omega$ sobre un recorrido de 2 k$\Omega$. La calibración manual redujo el error de magnitud del compensador respecto del modelo de 26,7 a 10,7 dB, pero el ajuste fino resulta difícil con el elemento actual. Se propone mantener la resistencia fija de la rama y utilizar un potenciómetro de menor recorrido y varias vueltas, dimensionado para cubrir la dispersión esperada de los componentes.

La precisión de ese ajuste adquiere mayor importancia en un arreglo de varios receptores. Diferencias entre componentes o entre posiciones de los potenciómetros modifican la magnitud y la fase de cada canal; la diferencia de fase entre circuitos se incorpora entonces a la que el método utiliza para estimar la propagación. El ajuste multivuelta facilitará igualar las respuestas, que deberán comprobarse experimentalmente entre placas.

La ruta PGA--ADC quedó identificada sobre 61 puntos, entre 0,21 Hz y 1,11 kHz, con coherencia mediana de 0,9946. El modelo identificado reproduce esos datos con residuos de 0,200 dB y 1,16°; estos valores describen el ajuste a la medición, no el error respecto de la respuesta nominal.

\subsection{Prueba controlada de adquisición}
\label{sec:prueba-controlada}

El segundo nivel verifica que el sistema adquiere de forma coordinada antes de exigirle un resultado geofísico. El nodo generador se trata como un esclavo más y comparte el arranque con el nodo receptor, de modo que cada respuesta queda referida al instante de su impacto. Para cuantificar esa referencia se midió, dentro de cada posición, la dispersión entre los instantes de disparo de los golpes individuales una vez alineados a nivel de submuestra, obteniéndose un valor del orden de 1,46 ms. Conviene precisar qué describe y qué no describe ese número. Es el residuo de repetibilidad entre golpes sucesivos de un mismo nodo, y agrupa el jitter mecánico del impacto con el del circuito de disparo. No mide, en cambio, ni la deriva entre jornadas de adquisición ni la incertidumbre de la posición del receptor, que son fuentes de error independientes y se tratan como parte del problema de geometría y repetibilidad de la campaña en las secciones~\ref{sec:campana} y~\ref{sec:limitaciones}.

Para $M$ golpes alineados, la traza promedio $\bar x[n]$ se obtiene de los registros $x_m[n]$ y sus desplazamientos estimados $\delta_m$:
\begin{equation}
\bar x[n]=\frac{1}{M}\sum_{m=1}^{M}x_m[n+\delta_m].
\label{eq:apilamiento}
\end{equation}
Los desplazamientos pueden ser fraccionarios y requerir interpolación. Si la señal se repite y el ruido es independiente entre golpes, su desviación estándar se reduce idealmente como $\sigma_{\bar n}=\sigma_n/\sqrt M$ \revsolution{\citep{Smith1997}}. Se registraron entre 20 y 37 golpes por posición; el apilamiento no elimina las variaciones de posición o de acoplamiento entre registros.

La capa de coordinación entre nodos sí fue probada de forma independiente. En laboratorio se operaron tres nodos simultáneos, uno de clase HAMMER y dos de clase GEO, que ejecutaron correctamente el arranque conjunto y el reparto de la orden de captura. Este ensayo se realizó sobre los módulos de radio sin el acondicionamiento analógico ni el PSoC asociados, de modo que demuestra que el protocolo de coordinación escala más allá de dos nodos y opera como fue diseñado, pero no aporta una cifra de precisión: no hubo conversión sincrónica de señal ni registros exportados que permitan comparar marcas temporales.

Conviene por eso separar cuatro magnitudes que resulta fácil confundir. La primera es la alineación de golpes recién mencionada, que corresponde al procesamiento de un solo nodo. La segunda es la operación funcional de la coordinación multinodo, demostrada con los tres nodos de laboratorio. La tercera es la sincronización digital medida entre nodos distintos, que exige registros simultáneos y una comparación de marcas temporales entre placas completas, y que no pudo realizarse por no disponerse de un conjunto de nodos GEO íntegros. La cuarta es la verificación instrumental del arranque coordinado mediante osciloscopio, que llegó a observarse durante el desarrollo pero cuyos registros no fueron preservados y que por lo tanto no constituye evidencia reproducible. La tolerancia de unos 670 µs derivada en la sección~\ref{sec:requerimientos} permanece, en consecuencia, como requerimiento de diseño cuya capa funcional está demostrada y cuya verificación metrológica está pendiente.

\subsection{Campaña de campo}
\label{sec:campana}

La campaña se realizó junto a una cancha de fútbol del predio de la Universidad Católica, sede Santa Librada, en Asunción. El punto de impacto permaneció fijo a varios metros del extremo del campo de fútbol y solamente se desplazó el nodo receptor, siguiendo una línea aproximadamente paralela a los laterales de la cancha. No se conservaron coordenadas ni azimut de los extremos del tendido, limitación que se retoma en la sección~\ref{sec:limitaciones}.

Se relevaron 21 posiciones entre 10 y 50 m de offset, con paso nominal de 2 m y una apertura total de 40 m, valores que corresponden a los límites de aliasing y de apertura establecidos en la sección~\ref{sec:requerimientos}. Del conjunto de 607 registros candidatos, 598 fueron incorporados al conjunto procesado tras la auditoría de calidad, con una mediana cercana a 30 golpes por posición. El arreglo se construyó desplazando un único receptor entre posiciones debido al número de transductores disponibles.

\begin{figure*}[!tb]
\centering
\includegraphics[width=\textwidth,keepaspectratio]{f_geometria_campanas.png}
\caption{Geometría de la campaña: 21 posiciones entre 10 y 50 m de offset con paso nominal de 2 m, lo que da una apertura total de 40 m. El arreglo se construyó desplazando un único receptor entre posiciones sucesivas. \revsolution{Fuente: elaboración propia a partir de los registros del proyecto.}}
\label{fig:geometria}
\end{figure*}

Esa restricción define con precisión qué valida la campaña. El conjunto permite construir un \textbf{registro multicanal sintetizado por posiciones con fuente activa repetible}, y sobre él evaluar la adquisición real, formar una imagen de dispersión e invertir un perfil preliminar. No demuestra, en cambio, la operación de 21 canales físicos simultáneos, ni la sincronización entre nodos, ni la igualdad de magnitud y fase entre placas distintas. La expresión «adquisición simultánea de 21 canales» sería incorrecta y no se emplea en ningún punto de este trabajo. Conviene señalar que el uso de un solo receptor no carece de antecedentes formales: Lin y Ashlock demostraron que la simulación multicanal con un receptor puede producir imágenes de dispersión comparables a las de MASW bajo condiciones controladas, si bien su protocolo fija el receptor y desplaza la fuente, es decir, exactamente lo inverso a lo realizado aquí \citep{LinAshlock2016}.

La ventana de offsets utilizada para el resultado final no se decidió durante la adquisición sino después, sobre el conjunto completo, y ese orden debe declararse para evitar la apariencia de una selección oportunista. El análisis posterior evaluó 26 subarreglos contiguos aplicando el mismo criterio a todos ellos, con el resultado de que las ventanas medio-lejanas superan tanto al arreglo completo como a las cercanas. La interpretación es compatible con una menor contaminación por campo cercano y por modos superiores al retirar los offsets pequeños, y con la pérdida de relación señal a ruido por atenuación en los más lejanos, aunque esta explicación no fue aislada mediante un ensayo independiente.

\begin{table*}[!ht]
\centering
\footnotesize
\caption{Barrido de subarreglos contiguos. \revsolution{Fuente: elaboración propia.}}\label{tab:subarreglos}
{\begin{tabularx}{0.92\textwidth}{YYYY}
\toprule
Ventana de offsets & Desajuste & Soporte de profundidad & Situación \\
\midrule
10–50 m (completo) & 2,28 \% & 8,81 m & Referencia \\
18–46 m & 1,85 \% & 10,06 m & Finalista (alcanza 11,83 m con otra convención) \\
22–42 m & 1,36 \% & 10,75 m & Base del modelo presentado \\
30–50 m & 16,5 \% & — & Descartada por desajuste \\
22–38 m & — & — & Descartada: fracción excesiva de la curva sobre el límite de aliasing \\
\bottomrule
\end{tabularx}}
\end{table*}

\section{Resultados preliminares}
\label{sec:resultados}

Los resultados se presentan en los mismos tres niveles de la sección anterior, distinguiendo en cada uno lo demostrado de lo interpretado.

\subsection{Resultados de la caracterización del instrumento}
\label{sec:resultado-instrumento}

La caracterización eléctrica cubre frecuencias inferiores a la resonancia del geófono y permite estudiar la respuesta del compensador en la banda que se busca recuperar. La transferencia medida, combinada con el modelo del sensor, proporciona una base para ajustar y evaluar esa extensión. El ensayo identificó también una mejora concreta del circuito: aumentar la resolución del ajuste del potenciómetro. La banda de campo se evalúa por separado, porque depende además de la excitación, del acoplamiento y del terreno.

\subsection{Resultados de la adquisición de campo}
\label{sec:resultado-campo}

La distribución de energía y de relación señal a ruido por bandas caracteriza el desempeño conjunto de la fuente, el terreno y el instrumento. Se calculó sobre 265 golpes repartidos en quince de las distancias relevadas, que es el subconjunto que superó el control de calidad exigido por esta métrica.

La banda de 10 a 50 Hz concentra la mayor parte de la energía registrada y presenta una coherencia mediana entre posiciones adyacentes de 0,732. Esa banda resume las métricas globales por frecuencia, pero no constituye un corte inferior absoluto: el procesamiento de dispersión utilizado conserva una curva desde aproximadamente 8 Hz. Se dispone así de resultados por debajo de los 10 Hz nominales del geófono, aunque con menor margen que en la banda de mayor energía.

El experimento permite identificar la ampliación de la apertura efectiva como un requerimiento para investigar frecuencias menores y capas más profundas. En el subarreglo de 22 a 42 m, cuya apertura es de 20 m, el punto de 8 Hz tiene una velocidad de fase de 172 m/s y una longitud de onda de 21,5 m, ya comparable con la extensión del arreglo. Su inclusión corresponde al criterio de análisis que admite longitudes de onda de hasta 1,5 veces la apertura. Para seguir la rama hacia frecuencias menores, con velocidades de fase similares o mayores, se necesitan longitudes de onda más grandes y un arreglo de mayor extensión que sostenga su interpretación.

La ampliación debe acompañarse de señal coherente en la banda de interés. Un tendido más extenso aporta soporte espacial, pero su aprovechamiento depende también de la energía de la fuente, el acoplamiento y la respuesta sensor--acondicionamiento. Los 8 Hz se conservan como extremo inferior de la curva utilizada, sin interpretarlos como el corte físico del instrumento.

\begin{figure*}[!tb]
\centering
\includegraphics[width=0.80\textwidth,keepaspectratio]{contenido_baja_frecuencia.png}
\caption{Contenido espectral del canal de geófono en campo. Arriba: densidad espectral de la señal por distancia fuente--receptor. Abajo: relación señal a ruido contra el ruido previo al disparo. El rótulo de 50 m corresponde a una meta histórica, retirada del proyecto. La curva de dispersión utilizada parte de 8 Hz; la mayor energía se concentra entre 10 y 50 Hz. \revsolution{Fuente: elaboración propia a partir de los registros del proyecto.}}
\label{fig:contenidobaja}
\end{figure*}

\subsection{Resultados geofísicos preliminares}
\label{sec:resultado-geofisico}

\subsubsection{Curva de dispersión y alcance}

La limitación geométrica para extender la interpretación hacia frecuencias menores está dada por la apertura efectiva del subarreglo, que condiciona la estimación de longitudes de onda grandes. Para estudiar capas más profundas se requiere una mayor apertura, acompañada de señal coherente en la banda de interés.

\begin{figure*}[!tb]
\centering
\includegraphics[width=\textwidth,keepaspectratio]{f_imagenes_dispersion.png}
\caption{Imágenes de dispersión en el plano frecuencia--velocidad de fase, normalizadas por frecuencia, de modo que los colores no comparan energía entre bandas. La línea dibujada es el límite de apertura del criterio de análisis adoptado. \revsolution{Fuente: elaboración propia a partir de los registros del proyecto.}}
\label{fig:dispersionimagenes}
\end{figure*}

\begin{figure*}[!tb]
\centering
\includegraphics[width=\textwidth,keepaspectratio]{f_curvas_dispersion.png}
\caption{Curvas de dispersión extraídas y su longitud de onda asociada, con la banda de incertidumbre obtenida por remuestreo. La curva empleada para el modelo presentado es la del subarreglo de 22 a 42 m. \revsolution{Fuente: elaboración propia a partir de los registros del proyecto.}}
\label{fig:dispersioncurvas}
\end{figure*}

El análisis y la selección de subconjuntos de las mediciones permitieron identificar longitudes de onda de 20 a 22 m. Para el punto de 8 Hz y 172 m/s del subarreglo de 22 a 42 m,
\begin{equation}
\lambda_R=\frac{172}{8}=21{,}5\ \mathrm m,
\qquad z_{\mathrm{sop}}\simeq\frac{\lambda_R}{2}=10{,}75\ \mathrm m.
\label{eq:alcance-campo}
\end{equation}
La variable $z_{\mathrm{sop}}$ representa el alcance estimado con esta convención, que se informa como aproximadamente \textbf{10 a 11 m}.

\subsubsection{Perfil e incertidumbre de la inversión}

Se adoptó un modelo de tres capas para la ventana de 22 a 42 m. La elección se basó en 24 remuestreos por \revsolution{\emph{bootstrap} \citep{Efron1979}} y 5 semillas del optimizador, con 120 inversiones por configuración. Con tres capas, prácticamente toda la variabilidad del perfil provino del remuestreo de los datos; al incorporar una cuarta, cerca de la mitad provino del optimizador. El dato disponible no permitió sostener esa complejidad adicional.

El perfil presenta una velocidad de corte somera del orden de 78 m/s, una interfaz efectiva próxima a 2,4 m y velocidades cercanas a 178 m/s por debajo, dentro del alcance investigado. La interfaz representa el cambio de rigidez que reproduce la curva observada.

La inversión supone una velocidad de corte creciente con la profundidad y un coeficiente de Poisson adoptado. En pruebas sintéticas, la primera restricción ocultó una inversión de velocidad conocida; al retirarla del análisis de campo, el perfil quedó insuficientemente determinado. Variar el coeficiente de Poisson afectó poco a la velocidad somera y aumentó la incertidumbre en profundidad. Por ello, el perfil constituye una interpretación compatible con los datos, sin identificar de forma única los contactos geológicos. \revsolution{La no unicidad de la inversión de ondas superficiales está documentada por \citet{Foti2009}; las pruebas anteriores describen cómo se manifestó en los datos de este trabajo.} El alcance de 10 a 11 m tampoco permite calcular $V_{S,30}$.

\subsection{Sensibilidad al procesamiento y a la selección de posiciones}
\label{sec:sensibilidad-procesamiento}

Cambiar la transformada modificó la curva de dispersión en 2,28 m/s de mediana; cambiar el algoritmo de extracción, en 1,89 m/s. Al retirar las posiciones de 32 y 40 m, la diferencia cuadrática media fue de 29,76 m/s. En las comparaciones realizadas, la selección de posiciones tuvo mayor efecto que la elección de la transformada o del algoritmo de extracción. Este resultado respalda priorizar el control de la geometría y la calidad de los registros en la siguiente campaña.

\subsection{Contraste con sondeos eléctricos del predio}
\label{sec:contraste-sondeos}

\revsolution{El informe hidrogeológico de julio de 2023, elaborado por Víctor González \& Asoc. S.S., aporta un contraste de plausibilidad \citep{EstudioHidrogeologico2023}.} Sus sondeos eléctricos verticales se ubican en el mismo predio, a distancias del orden de 120, 258 y 340 m de la línea de adquisición, que no pueden precisarse por no haberse registrado las coordenadas de los extremos. El sondeo más próximo presenta interfaces eléctricas a 1,00, 2,06 y 4,15 m, de modo que la interfaz efectiva de 2,4 m obtenida por inversión cae dentro de un intervalo en el que el sondeo eléctrico también detecta contactos. La coincidencia debe leerse como compatibilidad en el intervalo de 2 a 4 m y de ninguna manera como un acuerdo de centímetros, tanto por la distancia entre emplazamientos como porque ambas técnicas miden propiedades distintas: la resistividad no mide rigidez y la velocidad de corte no mide contenido de humedad. El contraste sirve para verificar que el resultado no es físicamente inverosímil, no para validarlo.

% Mantener el título de la síntesis unido a su tabla de ancho completo.
\begin{table*}[!ht]
\subsection{Síntesis: qué está demostrado y qué no}
\label{sec:sintesis}

\centering
\footnotesize
\caption{Síntesis del estado de cada afirmación. \revsolution{Fuente: elaboración propia.}}
\begin{tabularx}{0.92\textwidth}{YYY}
\toprule
Afirmación & Estado & Base \\
\midrule
La plataforma adquiere de forma conjunta una referencia de impacto y la respuesta del receptor & Demostrado & 598 registros aceptados con referencia temporal \\
Los datos se almacenan, organizan y procesan mediante un flujo reproducible & Demostrado & Manifiesto, servidor de ingesta y análisis repetible \\
La transferencia PGA--ADC se identifica entre 0,21 Hz y 1,11 kHz & Demostrado para una unidad & 61 puntos; residuos del ajuste de 0,200 dB y 1,16° \\
La energía de campo se concentra entre 10 y 50 Hz; la curva utilizada parte de unos 8 Hz & Medido; condicionado por el procesamiento & Métricas sobre 265 golpes en 15 distancias y curva de dispersión \\
Un arreglo sintetizado de 21 posiciones permite obtener imagen y curva de dispersión & Demostrado & Procesamiento de la campaña \\
El soporte de investigación es de aproximadamente 10 a 11 m & Preliminar & Longitud de onda soportada y convención $z\approx\lambda/2$ \\
El perfil de tres capas describe el sitio & Preliminar & 120 inversiones por configuración; sujeto a monotonía y a $\nu$ supuesto \\
La interfaz de 2,4 m corresponde a un contacto geológico & No afirmado & Es una interfaz efectiva del modelo \\
Adquisición simultánea con varios nodos GEO completos & No demostrado & Un solo nodo GEO íntegro operativo por registro \\
Coordinación de arranque con tres nodos simultáneos & Demostrado funcionalmente & Ensayo de laboratorio sobre los módulos de radio \\
Sincronización entre nodos dentro de unos 670 µs & No demostrado & Sin registros comparables entre placas completas \\
Ruido del instrumento & No medido & Caracterización pendiente \\
Extensión hacia capas más profundas y cálculo de $V_{S,30}$ & Trabajo futuro & El perfil actual tiene soporte hasta unos 10--11 m \\
\bottomrule
\end{tabularx}
\end{table*}

\section{Limitaciones y trabajos futuros}
\label{sec:limitaciones}

\textbf{Ruido del instrumento.} Se propone caracterizar el ruido para las configuraciones de ganancia y rango utilizadas.

\textbf{Operación multicanal.} Se completará la revisión de placa y la operación conjunta de varios nodos GEO, verificando la dispersión de ganancia y fase y la sincronización respecto de la tolerancia de unos 670 $\mu$s.

\textbf{Geometría y protocolo.} Se plantea registrar coordenadas y azimut, mantener una referencia de impacto y ensayar una mayor apertura efectiva para extender la caracterización por debajo de los 8 Hz actuales hacia capas más profundas.

\textbf{Fuente y extensión de la banda de excitación.} Se propone ensayar una fuente de impactos periódicos cuyo período pueda variarse. Para un período $T$, la frecuencia fundamental $f_p$ y sus armónicos $f_k$ cumplen
\begin{equation}
f_p=\frac{1}{T},\qquad f_k=\frac{k}{T},
\quad k=1,2,\ldots
\label{eq:fuente-periodica}
\end{equation}
donde $k$ es el orden del armónico \revsolution{\citep[cap.~13]{Smith1997}}. Al variar el período entre ensayos se desplazan la fundamental y sus armónicos. Se evaluará si esta excitación permite obtener registros coherentes a frecuencias menores.
